\documentclass[11pt]{article}
\usepackage[margin=1in]{geometry}
\usepackage{amsmath,amssymb}
\usepackage{booktabs}
\usepackage{hyperref}
\usepackage{xcolor}
\usepackage{enumitem}
\usepackage{titlesec}
\usepackage{graphicx}

\titleformat{\section}{\large\bfseries}{}{0em}{}[\titlerule]
\titleformat{\subsection}{\normalsize\bfseries}{}{0em}{}

\title{\textbf{Project Proposal II}\\[0.4em]
\large\textit{State-Aware Agent Unlearning for Toxicity Backflow in Social Thread Simulation}\\[0.3em]
\normalsize Target Venue: NeurIPS 2025}
\author{Anonymous}
\date{\today}

\begin{document}
\maketitle

% ─────────────────────────────────────────────────────────────
\section{Motivation and Problem Statement}
% ─────────────────────────────────────────────────────────────

Large language models deployed as conversational agents do not generate responses in
isolation. Each message they produce becomes part of the evolving context that conditions
all future responses---by themselves and by other agents in the same thread. This creates
a \textbf{stateful toxicity propagation} problem: a single toxic utterance early in a
thread can contaminate downstream responses even when later agents are individually
well-aligned, because the toxic content persists in the shared conversational state.

Existing mitigations are insufficient. Parameter-level LLM unlearning removes unsafe
behaviors from model weights but cannot prevent a fully-aligned model from being
re-conditioned on toxic content that re-enters through the conversation transcript.
Output moderation filters treat toxicity as a \emph{local output constraint}, but the
dominant failure mode in thread-based interaction is \textbf{contextual backflow}: toxic
content persists in the conditioning state and shapes downstream generations even after
its explicit spans are suppressed.

We study how to achieve \textbf{backflow-robust agent unlearning}: removing toxic
influence from an agent's behavior not only at the point of generation, but under
repeated re-exposure through the evolving conversational state.

% ─────────────────────────────────────────────────────────────
\section{Formal Problem Formulation}
% ─────────────────────────────────────────────────────────────

Given a human seed post $s$, we simulate a conversation as a directed graph $G=(V,E)$
where each node $v \in V$ is a message generated by an agent $A_i$, and each edge
$(u \to v)$ indicates that $v$ is produced conditioned on $u$.

We define the focal agent $A_1 \in \{\texttt{neutral}, \texttt{toxic}, \texttt{removed}\}$
and hold all other components fixed. The key quantities are:
\begin{align}
\mu(G) &= \frac{1}{|V|}\sum_{v \in V} \mathrm{tox}(v), \\
\Delta\mu(s) &= \mu(G^{(\texttt{toxic})}) - \mu(G^{(\texttt{neutral})}),
\end{align}
where $\mathrm{tox}(v) = f_{\text{tox}}(x_v) \in [0,1]$ is a continuous toxicity score.

\paragraph{Goal.} Given an agent embedded in an evolving discussion, learn an
unlearning intervention such that the agent does not (i)~produce toxic continuations,
nor (ii)~amplify toxicity in the broader thread, even when the environment continues
to surface toxic context through the conversational state.

% ─────────────────────────────────────────────────────────────
\section{Preliminary Results}
% ─────────────────────────────────────────────────────────────

\subsection{Experimental Setup}

We simulate Reddit-like conversation threads of the form $a \to b \to c \to d$ (4 turns),
where $A_1$ authors the first message $a$ under either a \texttt{neutral} or \texttt{toxic}
system prompt. All subsequent agents ($A_2, A_3, A_4$) are identical across conditions,
with no explicit toxicity instruction. Toxicity is scored using the Detoxify model across
six dimensions: toxicity, severe toxicity, identity attack, insult, obscene, and threat.
Reddit seed posts are used to provide realistic conversational grounding.

\subsection{Observed Propagation Signal}

Figure~\ref{fig:propagation} shows the mean Detoxify scores across turns 2--4
under neutral (blue) and toxic (orange) $A_1$ conditions.

\begin{figure}[h]
\centering
% Placeholder for the actual figure -- insert Screenshot_2026-03-26_at_6_32_49_PM.png
\fbox{\parbox{0.92\textwidth}{\centering\vspace{2em}
\textit{[Figure: Mean Detoxify scores by dimension (identity\_attack, insult, obscene,
severe\_toxicity, threat, toxicity) across turns 2--4 under neutral vs.\ toxic $A_1$.
Toxic condition consistently higher across all six dimensions.]}
\vspace{2em}}}
\caption{Toxicity propagation by dimension. Toxic $A_1$ elevates all six Detoxify
dimensions in downstream responses compared to a neutral $A_1$, confirming stateful
influence propagation.}
\label{fig:propagation}
\end{figure}

\paragraph{Key observations from preliminary results.}

\begin{itemize}[noitemsep]
  \item The toxic condition (orange) is \textbf{consistently higher than neutral (blue)
  across all six dimensions}, confirming that toxicity functions as a form of stateful
  influence in closed-loop agentic interaction.
  \item The strongest signal appears in the \textbf{toxicity} dimension (peak $\approx 0.06$
  vs.\ $\approx 0.004$ for neutral) and the \textbf{insult} dimension (peak $\approx 0.017$
  vs.\ $\approx 0.0001$), representing roughly a 15$\times$ and 170$\times$ elevation
  respectively.
  \item The \textbf{obscene} dimension shows a monotonically increasing trend under the
  toxic condition across turns 2--4, suggesting cumulative drift rather than immediate
  saturation---a key feature of stateful backflow.
  \item Absolute values are modest overall, consistent with safety-aligned base models
  resisting overtly toxic generation even under toxic context. This motivates the
  signal-amplification strategies described in Section~\ref{sec:experiments}.
\end{itemize}

\paragraph{Gap and direction.}
The phenomenon is real but the effect size is small under the current setup. The main
causes are: (i)~modern safety-tuned LLMs resist generating explicitly toxic content even
when primed; (ii)~a 4-turn chain is too short to accumulate significant drift; and
(iii)~the current toxicity dose uses a single binary \texttt{toxic}/\texttt{neutral}
condition without a graded dose parameter. All three are addressable.

% ─────────────────────────────────────────────────────────────
\section{Proposed Method: State-Aware Read/Write Unlearning}
% ─────────────────────────────────────────────────────────────

\subsection{Core Framework}

We treat agent unlearning as \textbf{state control} with two intervention points:

\paragraph{Read-side sanitization (context filtering).}
Before passing context $\mathcal{C}(v)$ to an agent, we transform it with a sanitizer $S$:
\[
\widetilde{\mathcal{C}}(v) = S(\mathcal{C}(v)).
\]
Instantiations:
\begin{itemize}[noitemsep]
  \item \textbf{Redaction:} remove message spans whose toxicity exceeds threshold $\tau$.
  \item \textbf{Constrained summarization:} summarize visible history while forbidding toxic spans.
  \item \textbf{Soft filtering:} mask high-toxicity segments while retaining structural cues.
\end{itemize}

\paragraph{Write-side gating (state update control).}
After generating message $x_v$, control what is appended to the transcript:
\[
\texttt{state} \leftarrow \texttt{state} \cup W(x_v,\, \mathrm{tox}(v)),
\]
where $W$ stores either the raw message (baseline) or a sanitized version / refusal stub
if $\mathrm{tox}(v) > \tau$. This prevents self-reinfection where toxic outputs become
future conditioning context.

\subsection{DPO Extension (Parameter-Level)}

As an optional extension---and to connect to the preference unlearning literature---we
apply DPO at the agent level:
\begin{equation}
\mathcal{L}_{\text{pref}}
= -\log \sigma\!\Big(\beta\big[\log \pi_\theta(m_t^+ \mid \widetilde{\mathcal{C}}(v)) - \log \pi_\theta(m_t^- \mid \mathcal{C}(v))\big]\Big),
\end{equation}
where $m_t^+$ is the response conditioned on sanitized context and $m_t^-$ is the
response conditioned on raw toxic context.
The key empirical claim is that \textbf{state-level control is necessary even when
parameter-level DPO is applied}: a well-trained policy can still be compromised by
toxic context re-entering through the transcript state.

% ─────────────────────────────────────────────────────────────
\section{Experimental Plan}
\label{sec:experiments}
% ─────────────────────────────────────────────────────────────

\subsection{Week 1: Amplify the Signal}

\begin{enumerate}[noitemsep]
  \item \textbf{Extend chain length} from 4 to 8 turns to allow cumulative drift to accumulate.
  \item \textbf{Graded toxicity dose} $\alpha \in \{0, 0.25, 0.5, 0.75, 1.0\}$ via stochastic
  prompt mixing between a neutral system prompt $P_N$ and a toxic prompt $P_T$:
  with probability $\alpha$ use $P_T$, else $P_N$.
  \item \textbf{Subtler toxicity dimensions:} reframe toxicity as condescension, sarcasm,
  or identity-coded dismissiveness---forms that bypass safety filters while remaining
  measurable by Detoxify. Use adversarial Reddit seed posts from toxic subreddits.
  \item \textbf{Target:} a mean $\Delta\mu \geq 0.05$ on the toxicity dimension across
  seeds before proceeding to mitigation experiments.
\end{enumerate}

\subsection{Week 2: Core Ablation Table}

We compare four intervention conditions using the paired counterfactual protocol:
\begin{center}
\begin{tabular}{lcc}
\toprule
\textbf{Method} & \textbf{Toxicity $\downarrow$} & \textbf{Coherence $\uparrow$} \\
\midrule
None (baseline) & -- & -- \\
Read-only (context sanitization) & & \\
Write-only (state gating) & & \\
Joint (read + write) & & \\
\bottomrule
\end{tabular}
\end{center}

All comparisons use the same seed, topology, and agent assignment to isolate the effect
of the intervention.

\subsection{Metrics}

\paragraph{Global vibe shift.}
\[
\mu(G)=\frac{1}{|V|}\sum_{v\in V}\mathrm{tox}(v), \qquad
\Delta\mu(s)=\mu(G^{(\texttt{toxic})})-\mu(G^{(\texttt{neutral})}).
\]

\paragraph{Contagion by distance.}
\[
g(k)=\mathbb{E}[\mathrm{tox}(v)\mid d(v,A_1)=k],\quad k=1,2,\dots,K.
\]
This directly visualizes propagation radius and decay under mitigation.

\paragraph{Temporal dynamics.}
\begin{itemize}[noitemsep]
  \item \textbf{Time-to-first-toxic (TTFT):} earliest turn $t$ where $\mathrm{tox}(v_{(t)}) > \tau$.
  \item \textbf{Area Under Toxicity Curve (AUTC):} $\sum_{t=1}^{|V|}\mu_t$, capturing both severity and persistence.
\end{itemize}

\paragraph{Multi-run statistics.}
Each scenario is run with $R = 5$ random seeds; we report mean and worst-case
($\mu_{\max}$) to capture leakage under stochastic sampling.

\subsection{Week 3: Topology and Exposure Ablations}

\begin{itemize}[noitemsep]
  \item \textbf{Topology ablation:} evaluate chain, balanced tree (depth 3, branching 2),
  cross-linked DAG, and high-branching tree. Test whether toxicity amplification
  increases with branching factor.
  \item \textbf{Exposure ablation:} vary conditioning set $\mathcal{C}(v)$ across
  parent-only, path-to-root, thread-local, and full-visible regimes to isolate
  whether propagation is driven by transcript exposure (backflow) vs.\ intrinsic agent bias.
  \item \textbf{Population dose:} vary $\rho$ (fraction of toxic agents) to test
  community-level amplification independently of individual dose $\alpha$.
\end{itemize}

\subsection{Week 4: Writing and NeurIPS Checklist}

Complete limitations section, reproducibility details (model versions, compute estimates,
random seed protocols), ethical considerations (Reddit data usage, toxicity annotation),
and NeurIPS checklist items.

% ─────────────────────────────────────────────────────────────
\section{Contributions and Novelty}
% ─────────────────────────────────────────────────────────────

\begin{enumerate}[noitemsep]
  \item \textbf{Threat model:} We distinguish local toxicity (agent output) from system
  toxicity (downstream propagation through evolving state), and define relapse under
  repeated environmental exposure.
  \item \textbf{State-centric unlearning framework:} Read-side context sanitization and
  write-side state gating as model-agnostic mechanisms for backflow-robust agent unlearning.
  \item \textbf{Controlled evaluation protocol:} Paired counterfactual simulation with
  fixed topology, ordering, and randomness---isolating the causal effect of a single
  agent's toxicity on downstream thread quality.
  \item \textbf{Empirical finding:} Preliminary results confirm that toxicity propagates
  statefully across all six Detoxify dimensions; the obscene and toxicity dimensions show
  the strongest and most sustained drift under toxic priming.
\end{enumerate}

% ─────────────────────────────────────────────────────────────
\section{Risks and Mitigations}
% ─────────────────────────────────────────────────────────────

\begin{tabular}{p{5.5cm}p{8.5cm}}
\toprule
\textbf{Risk} & \textbf{Mitigation} \\
\midrule
Effect size too small for strong claims & Extend chain length, use graded dose, target subtler toxicity dimensions. \\
Safety-tuned models resist toxic generation & Use adversarial seed posts; consider less-safety-tuned model for simulation (not for mitigation). \\
Read/write controls degrade coherence & Report output coherence alongside toxicity; use constrained summarization rather than hard redaction. \\
Crowded toxicity space at NeurIPS & Emphasize the \emph{agentic/stateful} angle; position as unlearning paper, not toxicity detection. \\
\bottomrule
\end{tabular}

% ─────────────────────────────────────────────────────────────
\section{Relationship Between the Two Projects}
% ─────────────────────────────────────────────────────────────

This proposal and the companion privacy leakage proposal (Proposal I) share a common
theoretical framework: both treat unsafe agentic behavior as a \emph{state propagation
problem} amenable to read/write control and preference-based unlearning. The two projects
can be developed in parallel and potentially unified into a single paper that demonstrates
the generality of the state-centric unlearning framework across two distinct harm modalities
(toxicity and privacy). If submitted separately, each contributes a self-contained
empirical story. If combined, the unifying framework becomes the primary contribution.

\end{document}
